\section{Empirical Measurement}
\label{sec:empirical}

\subsection{Subjects and method}

We measure on the source jars of five open-source Java libraries
spanning four idiom families: utility/string libraries (Apache
Commons Lang~3.14.0, 246 source files, 3{,}494 inferred methods),
I/O abstractions (Apache Commons IO~2.13.0, 245 files, 2{,}111
methods), numerical algorithms (Apache Commons Math~3.6.1, 990
files, 6{,}645 methods), functional-style combinators (jOOL~0.9.14,
130 files, 3{,}691 methods), and persistent collections (Vavr~0.10.4,
101 files, 5{,}111 methods). Combined, the corpus is 1{,}712 source
files and 21{,}052 methods with a non-empty isolated baseline spec,
giving the measurement a breadth that captures both the imperative
utility-class idiom emphasised by Commons Lang/IO and the
higher-order generic style emphasised by jOOL/Vavr. Commons Lang and
IO are also used by the predecessor work~\cite{edmonds2026inference,
edmonds2026embedding}, which keeps cross-paper comparisons free of
corpus-difference confounds for those two libraries.

For each library, our harness:

\begin{enumerate}\setlength\itemsep{0.2em}
  \item Parses every \texttt{.java} entry in the sources jar.
  \item Builds the call graph
        (\texttt{CallGraphBuilder.\allowbreak{}buildFromCompilationUnits}).
  \item Runs the existing
        \texttt{Method\allowbreak{}Specification\allowbreak{}Inferrer.\allowbreak{}infer\allowbreak{}Specification}
        on every method declaration. The inferrer's internal
        \texttt{InterproceduralAnalyzer} runs as part of this pass,
        propagating callee preconditions in the standard single-pass
        manner. The resulting per-method precondition and
        postcondition counts are snapshotted as the
        \emph{isolated baseline}.
  \item Invokes
        \texttt{CompositionalAnalyzer.\allowbreak{}refineAll()} on
        the populated cache, which mutates each method's spec in
        place.
  \item For each method, diffs the post-refinement spec against the
        baseline and counts newly-added clauses. Each new
        precondition is categorised by syntactic shape (null check,
        range/bound, branch-conditional implication, disjunctive,
        other) for the RQ2 distribution.
\end{enumerate}

All measurements are deterministic given the inputs; we report each
exactly once. The harness, the per-library output, and the
configuration knobs are committed to the replication package.

\subsection{Results}

Table~\ref{tab:headline} reports the headline measurements across
all five libraries.

\begin{table*}[ht]
\centering
\small
\caption{Effect of compositional refinement over the isolated single-pass baseline across five libraries. \emph{Refined} counts methods whose spec gained at least one new clause; percentages are of \emph{methods in cache}. The pass extends the spec on 23--43\% of inferred methods across every library, producing 45{,}740 new precondition clauses in total.}
\label{tab:headline}
\begin{tabular}{lrrrrr}
\toprule
& \textbf{Lang} & \textbf{IO} & \textbf{Math} & \textbf{jOOL} & \textbf{Vavr} \\
\midrule
Methods in cache             & 3{,}494 & 2{,}111 & 6{,}645 & 3{,}691 & 5{,}111 \\
\midrule
Methods refined              & 1{,}513 & 750     & 2{,}417 & 862     & 1{,}728 \\
\hspace{1em}\emph{as \% of cache} & 43.30\% & 35.53\% & 36.37\% & 23.35\% & 33.81\% \\
Methods with new precondition & 1{,}137 & 514    & 1{,}883 & 768     & 1{,}223 \\
Methods with new postcondition & 1{,}246 & 655   & 1{,}894 & 447     & 1{,}260 \\
\midrule
New precondition clauses     & 5{,}755  & 1{,}983 & 27{,}240 & 7{,}309 & 3{,}453 \\
\hspace{1em}\emph{avg per refined method} & 5.06 & 3.86 & 14.47 & 9.52 & 2.82 \\
\hspace{1em}\emph{max on one method} & 216   & 41 & 710 & 208 & 177 \\
New postcondition clauses    & 5{,}166  & 2{,}086 & 11{,}613 & 3{,}156 & 2{,}874 \\
\bottomrule
\end{tabular}
\end{table*}

\textbf{RQ1.} The second pass strictly extends the isolated set on every library. The refinement adds at least one
clause to 23.35--43.30\% of cached methods (median 35.53\%, IQR
33.81--36.37\%). In total the pass adds 45{,}740 new precondition
clauses and 24{,}895 new postcondition clauses across the five
libraries.

The single-pass propagation is \emph{not} subsuming, and
the gap is consistent across idiom families: a utility library
(Commons Lang) shows the highest fraction of refined methods because
its delegating helpers expose many guarded call sites, while a
generic functional combinator library (Vavr) shows fewer but the
new clauses concentrate in disjunctive shapes from polymorphic
dispatch. The result holds even though all five libraries' isolated
pipelines already invoke
\texttt{InterproceduralAnalyzer}'s single-pass callee-precondition
propagation; the second pass discovers obligations the first pass
cannot express.

\subsubsection{Per-method refinement distribution}

Aggregate counts obscure the per-method dispersion: a few large
methods could in principle account for most of the new clauses. The
ratio of \emph{methods refined} to \emph{total new clauses} (mean
new preconditions per refined method) bounds the per-method
dispersion above: 5.06 on Commons Lang, 3.86 on Commons IO, 14.47
on Commons Math, 9.52 on jOOL, and 2.82 on Vavr. Even on the
heaviest library (Math), the per-refined-method mean is well below
the reported maxima (710 on Math, 216 on Lang, 208 on jOOL, 177 on
Vavr, 41 on IO), establishing that the maxima are deep-tail outliers
and the bulk of the population contributes a small number of clauses
each. The refinement is therefore a population-level effect on
every library, not an artefact of a handful of large methods.

\subsection{Shape distribution}

Table~\ref{tab:shapes} categorises the newly-added preconditions by
syntactic shape. The categorisation runs on each new clause
independently and is structurally exhaustive (every clause falls in
exactly one bucket; ties resolve in favour of structural shape
over content, so a clause containing \texttt{==>} is counted as
branch-conditional even if it also mentions \texttt{!= null}).

\begin{table*}[ht]
\centering
\small
\caption{Syntactic shape of the newly-added preconditions across the five libraries. Percentages are of each library's total new preconditions. Branch-conditional implications are the largest single shape on four libraries out of five; on jOOL, null checks dominate because of pervasive generic input validation. Disjunctive clauses, the polymorphic-dispatch signal, are present on every library and reach 9.8\% on Vavr.}
\label{tab:shapes}
\begin{tabular}{lrrrrrrrrrr}
\toprule
& \multicolumn{2}{c}{\textbf{Lang}} & \multicolumn{2}{c}{\textbf{IO}} & \multicolumn{2}{c}{\textbf{Math}} & \multicolumn{2}{c}{\textbf{jOOL}} & \multicolumn{2}{c}{\textbf{Vavr}} \\
\cmidrule(lr){2-3}\cmidrule(lr){4-5}\cmidrule(lr){6-7}\cmidrule(lr){8-9}\cmidrule(lr){10-11}
\textbf{Shape} & count & \% & count & \% & count & \% & count & \% & count & \% \\
\midrule
Branch-conditional (\texttt{==>})       & 3{,}329 & 57.8\% & 910   & 45.9\% & 18{,}273 & 67.1\% & 1{,}899 & 26.0\% & 1{,}939 & 56.2\% \\
Null check (\texttt{!= null}/\texttt{== null}) & 1{,}152 & 20.0\% & 652   & 32.9\% & 6{,}025  & 22.1\% & 5{,}103 & 69.8\% & 1{,}030 & 29.8\% \\
Range / bound (\texttt{<,<=,>,>=})       & 926     & 16.1\% & 350   & 17.6\% & 1{,}658  & 6.1\%  & 23      & 0.3\%  & 123     & 3.6\%  \\
Disjunctive (\texttt{||})                & 260     & 4.5\%  & 36    & 1.8\%  & 1{,}002  & 3.7\%  & 274     & 3.7\%  & 338     & 9.8\%  \\
Other                                    & 88      & 1.5\%  & 35    & 1.8\%  & 282      & 1.0\%  & 10      & 0.1\%  & 23      & 0.7\%  \\
\midrule
Total new preconditions                  & 5{,}755 & 100\%  & 1{,}983 & 100\% & 27{,}240 & 100\%  & 7{,}309 & 100\%  & 3{,}453 & 100\%  \\
\bottomrule
\end{tabular}
\end{table*}

\textbf{RQ2.} Branch-conditional implications and disjunctions --- the two shapes the single pass structurally cannot produce --- are present on every library and dominate on most.
Branch-conditional clauses are the largest single shape on four of
the five libraries, peaking at 67.1\% on Commons Math (where heavy
algorithmic code paths gate calls on numerical preconditions). The
single exception is jOOL, where null checks dominate (69.8\%)
because the higher-order functional combinators delegate to
user-supplied lambdas through a small set of pinch-point validation
sites; even there, branch-conditional shapes account for 26.0\%, a
substantial quarter.

Disjunctive clauses --- the polymorphic-dispatch signal --- are present
on every library and peak at 9.8\% on Vavr, the library with the
deepest interface inheritance chains (\texttt{Value} $\to$
\texttt{Traversable} $\to$ \texttt{Iterable} $\to \cdots$). Both
shapes confirm the structural claim: the single pass, which lacks
control-flow context and resolves to a single nominal callee, cannot
produce them.

Null and range/bound clauses (the remaining ~21\% on Lang/IO,
~28\% on Math, ~33\% on Vavr, ~70\% on jOOL) are obligations the
single pass \emph{would} produce when the call site is unguarded,
but the second pass adds because the call site is guarded by a
condition the single pass would not lift. The residual \emph{other}
(0.1--1.8\%) is the long-tail shape that includes arithmetic
comparisons not matched by the range/bound regex.

\subsection{A version-step instance: Commons Lang 3.12.0 \texorpdfstring{$\to$}{->} 3.14.0}
\label{sec:empirical-version}

The cross-method propagation we measure has an immediate consequence: a change to a library propagates into the inferred contracts of methods that call it. We ran the inferrer over two consecutive Commons Lang versions (3.12.0 and 3.14.0) and diffed the per-method inferred specifications (Table~\ref{tab:version-diff}), to show this propagation operates at scale on a real version step rather than only in principle.

\begin{table}[ht]
\centering
\small
\caption{Inferred-spec diff on Commons Lang 3.12.0 vs 3.14.0. The propagation mechanism this paper measures surfaces a tractable, classifiable candidate-breaking set on a real two-year version step.}
\label{tab:version-diff}
\begin{tabular}{lr}
\toprule
Methods with inferred spec, 3.12.0     & 3{,}216 \\
Methods with inferred spec, 3.14.0     & 3{,}494 \\
Methods present in both versions       & 3{,}188 \\
\midrule
Inferred spec unchanged                & 2{,}663 \\
Inferred spec \textbf{changed}         & \textbf{525} \\
\hspace{1em} precondition strengthened ($+$req)        & 144 \\
\hspace{1em} precondition weakened ($-$req)            & 87 \\
\hspace{1em} precondition mixed                        & 90 \\
\hspace{1em} postcondition strengthened ($+$ens)       & 78 \\
\hspace{1em} postcondition weakened ($-$ens)           & 48 \\
\hspace{1em} postcondition mixed                       & 178 \\
\midrule
Methods only in 3.12.0 (removed)       & 28 \\
Methods only in 3.14.0 (added)         & 306 \\
\bottomrule
\end{tabular}
\end{table}

Of the 3{,}188 methods present in both versions, 525 (16.5\%) have an
inferred specification that differs --- exactly the candidate-breaking set a
behavioural-compatibility checker would surface. 144 have a \emph{strengthened}
precondition (a clause the new version requires that the old one did not):
client code that called the old version may now violate the added requirement.
A spot check of these confirms a mix of real precondition tightenings
(\textit{e.g.}, added non-null guards on instance fields of
\texttt{CharRange}, \texttt{CharSet}, and the various \texttt{Builder}
classes) and inferrer noise. Validating each against the release notes is
left to a separate end-to-end study, but the candidate set is the right size
and the right shape for that validation to be tractable: it is the
mechanism's first practical output on a real version step.

\subsection{Runtime cost}

Table~\ref{tab:runtime} reports the refinement wall-clock per
library. The second pass completes in 2.2 seconds on Commons Lang,
0.7 seconds on Commons IO, and scales sub-linearly through 13.4
seconds on Commons Math (the largest library at 6{,}645 methods
across 990 source files). No SCC hit the 16-iteration fixpoint cap
on any library; the call graphs of all five corpora are acyclic
under the analyser's signature-matching scheme, so each method is
refined exactly once. The refinement step is therefore an addition
of seconds-to-tens-of-seconds on top of an inference pass that
already takes minutes on libraries of this size; it is a clearly
affordable enhancement.

\begin{table}[ht]
\centering
\small
\caption{Refinement wall-clock per library. The pass scales close to linearly with the number of SCCs visited; the per-method cost is dominated by the AST walk for call sites.}
\label{tab:runtime}
\begin{tabular}{lrrr}
\toprule
\textbf{Library} & \textbf{Methods} & \textbf{SCCs} & \textbf{Refine ms} \\
\midrule
commons-lang3-3.14.0   & 3{,}494 & 6{,}223  & 2{,}164.0 \\
commons-io-2.13.0      & 2{,}111 & 3{,}994  & 664.4     \\
commons-math3-3.6.1    & 6{,}645 & 13{,}342 & 13{,}413.7\\
jool-0.9.14            & 3{,}691 & 5{,}403  & 2{,}790.6 \\
vavr-0.10.4            & 5{,}111 & 7{,}610  & 15{,}039.4\\
\bottomrule
\end{tabular}
\end{table}

\textbf{RQ3.} The second pass is cheap and fixpoint iteration never fires at the libraries' scale. Total refinement
wall-clock across all five libraries is 34 seconds for 21{,}052
methods (1.62 ms/method on average). The runtime characteristic is
dominated by the AST walk per method
(\texttt{decl.findAll(MethodCallExpr.class)}), which is linear in
the body size; the lifting and substitution steps add constant
factors per call site.

